\documentclass[11pt]{article}
\usepackage[margin=0.82in]{geometry}
\usepackage{amsmath,amssymb,booktabs,graphicx,xurl,hyperref}
\usepackage[T1]{fontenc}
\newcommand{\NA}{\textemdash}
\title{Symmetry-sector follow-up of the three-mode autocorrelations}
\author{}
\date{7 September 2026}
\begin{document}
\maketitle

\section*{Scope and predeclared analysis}
This is an analysis-only follow-up using the saved correlations from 64
independent trajectories per bath condition.  No simulation was run.  The
original plateau gate is unchanged.  Before fitting, the early damped-fit
window was fixed to $t\in[0.05,1.00]$ for both observables and both bath
conditions.  The trajectory bootstrap uses 2,000 resamples.

A terminology correction is essential. Under endpoint exchange
$1\leftrightarrow3$, $\theta_1\leftrightarrow\theta_3$; hence
$\sin\theta_1$ alone is not exchange odd, and $\cos\theta_1$ and
$\cos\theta_3$ alone are not invariant. The saved data do not contain
$\sin\theta_3$, so they do not span exact exchange sectors. We therefore call
the requested five observables the \emph{even-candidate group} and
$\{\sin\theta_1,I_1-I_3\}$ the \emph{sign-changing pair}. Only $I_1-I_3$ is
strictly exchange odd. In the driven case the unequal baths break exchange
symmetry explicitly.

\section*{Even-candidate plateau test}
\begin{table}[h]
\centering\small
\begin{tabular}{lcccc}
\toprule
case & individually resolved & rate range & relative spread & common result\\
\midrule
$(2,8)$ & $I_2,\ I_1+I_3$ & $[-1.109,-0.993]$ & 11.1\% & no\\
$(5,5)$ & $\cos\theta_1,\ \cos\theta_3,\ I_2$ & $[-1.188,-0.907]$ & 24.8\% & no\\
\bottomrule
\end{tabular}
\end{table}

For the driven case, the two accepted rates overlap in time on
$[1.12,1.58]$ and their confidence intervals overlap, but the frozen common
gate requires at least three independently resolved observables. The other
three candidates have no accepted plateau, so no group rate is quoted. At
equal temperature three candidates pass individually and overlap in time on
$[1.00,1.68]$, but the spread exceeds 20\% and their 95\% intervals have empty
intersection: the largest lower endpoint is $-0.9665$, while the smallest
upper endpoint is $-1.0160$. Thus neither condition resolves a common slow
rate for the requested five-observable group.

\section*{Early sign-changing fits}
The fitted model is
\[
 C(t)=e^{\lambda_Rt}\{A\cos(\lambda_It)+B\sin(\lambda_It)\}.
\]
Frequency resolution requires a 95\% interval excluding zero,
$\Delta\mathrm{AIC}\le-10$ against a single real exponential, and at least
half a fitted cycle in the fixed window,
$\lambda_I\ge\pi/0.95=3.30694$.

\begin{table}[h]
\centering\scriptsize
\begin{tabular}{llrrrrr}
\toprule
case & $g$ & $\lambda_R$ (95\% CI) & $\lambda_I$ (95\% CI) & period (95\% CI) & $\Delta$AIC & verdict\\
\midrule
$(2,8)$ & $\sin\theta_1$ & $-4.220$ $[-4.270,-4.172]$ & $3.306$ $[3.112,3.495]$ & $1.901$ $[1.798,2.019]$ & $-617.8$ & unresolved\\
$(2,8)$ & $I_1-I_3$ & $-4.911$ $[-5.146,-4.694]$ & $5.356$ $[5.205,5.507]$ & $1.173$ $[1.141,1.207]$ & $-516.5$ & resolved\\
$(5,5)$ & $\sin\theta_1$ & $-5.006$ $[-5.069,-4.945]$ & $3.896$ $[3.714,4.066]$ & $1.613$ $[1.545,1.692]$ & $-466.4$ & resolved\\
$(5,5)$ & $I_1-I_3$ & $-4.285$ $[-4.427,-4.147]$ & $5.351$ $[5.200,5.499]$ & $1.174$ $[1.143,1.208]$ & $-986.8$ & resolved\\
\bottomrule
\end{tabular}
\end{table}

The driven $\sin\theta_1$ estimate misses the frozen half-cycle threshold by
$9.1\times10^{-4}$ and its bootstrap interval straddles that threshold. It is
therefore reported as unresolved, despite its strong AIC improvement. Every
fit accepted 2,000/2,000 bootstrap resamples.

\begin{figure}[h]
\centering
\includegraphics[width=\textwidth]{sector_followup/early_damped_sector_fits.pdf}
\caption{Saved mean correlations and the fixed-window damped and
real-exponential fits. Shading is the trajectory-level 95\% pointwise band;
vertical dotted lines mark the interpolated zero of the saved mean.}
\end{figure}

\section*{Zero-crossing consistency}
\begin{table}[h]
\centering\small
\begin{tabular}{llrrr}
\toprule
case & $g$ & mean zero & fitted zero (95\% CI) & first $3\,$SE-negative lag\\
\midrule
$(2,8)$ & $\sin\theta_1$ & 0.2726 & 0.2752 $[0.2726,0.2775]$ & 0.28\\
$(2,8)$ & $I_1-I_3$ & 0.4211 & 0.4189 $[0.4129,0.4254]$ & 0.44\\
$(5,5)$ & $\sin\theta_1$ & 0.2326 & 0.2356 $[0.2339,0.2372]$ & 0.24\\
$(5,5)$ & $I_1-I_3$ & 0.3891 & 0.3891 $[0.3852,0.3935]$ & 0.40\\
\bottomrule
\end{tabular}
\end{table}
The fitted phase predicts the observed first zero accurately in all four fits.
The first significantly negative lag is slightly later, as expected from its
$3\,$SE support requirement. A period alone cannot predict the first zero,
because the zero also depends on the fitted phase.

\section*{Sector comparison and claim boundary}
The two sign-changing observables do not agree on either $\lambda_R$ or
$\lambda_I$: their 95\% intervals are disjoint in both bath conditions. Their
real parts, approximately $-4.2$ to $-5.0$, are also far faster than the
individually resolved even-candidate rates near $-1$. Equal temperature does
not produce cleaner agreement between this requested pair. This is expected
once $\sin\theta_1$ is recognized as a mixture rather than an exact odd-sector
observable.

The defensible conclusion is that $I_1-I_3$ exhibits a well-resolved early
damped mode, and $\sin\theta_1$ contains a distinct sign-changing component.
The strong AIC preference over a \emph{single} real exponential and the
accurate zero crossings are consistent with oscillatory relaxation, but do not
alone prove a complex generator eigenpair: a signed sum of several real modes
can also cross zero. Exact sector identification would require saving
$\sin\theta_3$ and analysing the sums and differences under endpoint exchange.

Machine-readable results are under \path{sector_followup/results/}; the frozen
analysis specification and input hashes are in
\path{sector_followup/PROTOCOL.md}.
\end{document}
